\subsection{Docker Swarm}
Docker selber bietet als Cluster-Orchestrator Docker Swarm an (zureinfachheit absofort nur Swarm genannt) und ist so mit einer der einfachsten Wege um
vom einen Compose zu einen Cluster-Orchestrator zuwechseln, da weiterhin die compose.yaml verwendet wird.

Die compose.yaml Datei muss dafür mit dem Befehl docker stack deploy ausgeführt werden, dabei setzt Swarm aber vorraus, dass das Gerät zu einen Cluster gehöhrt bzw 
wie es in Swarm heißt, einen Docker-Swarm. Es würd aber nicht verlangt, das der Swarm schon aus mehreren Host besteht,
sondern kann mit dem Befehl docker swarm init, ein Swarm mit nur einen Host auf den Rechner erstellt werden \parencite[vgl. 119]{kofler2023docker}, 
daher kann Swarm auch nur auf einen Node verwendet werden, auch wenn es als Multi-Node ausgelegt ist.

Dennoch besteht ein Swarm immer aus zwei verschiedenen Arten von Nodes. 
Einerseits gibt es die Manager Nodes, die die Cluster-Verwaltung übernehmen, andererseits gibt es die Worker Nodes, auf denen die Container ausgeführt werden. 
Ein Manager Node kann aber auch zusätzlich als Worker Node verwendet werden, beispielsweise wenn der Swarm nur aus einem Node besteht.

Die Vorteile von Swarm im Vergleich zu Compose sind zum einen, dass Swarm direkt auf mehrere Nodes skaliert werden kann und dass die Replikate der Container nun hochskaliert werden können. 
Swarm sorgt mit seiner Cluster-Verwaltung außerdem dafür, dass die Anzahl der Replicas der Container immer gleich bleibt, auch wenn ein Container oder sogar ein ganzer Node ausfällt.
Dieser Zustand wird als Desired-State-Prinzip bezeichnet.

Obwohl Swarm einen einfachen Übergang zu einer Cluster-Orchestrierung bietet, ist es dennoch nicht die optimale Wahl als Cluster-Orchestrator.
Einerseits ist die Anwendung weiterhin an Docker-Techniken gebunden, andererseits ist Swarm nicht so stark verbreitet, weswegen es eine kleinere Community hat.
Dadurch ist Swarm weniger erweiterbar und flexibler als andere Cluster-Orchestratoren.
